\documentclass{xduugthesis}

\xdusetup{
    style = {
        cjk-font = fandol,
        latin-font = gyre
    },
    info = {
        title = {基于 AASIST 的中文语音克隆检测系统设计与实现},
        author = {作者姓名},
        department = {学院名称},
        major = {专业名称},
        supervisor = {指导教师},
        student-id = {学号},
        abstract = {chapters/abstract-zh.tex},
        keywords = {语音克隆检测, 语音反欺骗, AASIST, 图注意力网络, 深度伪造语音},
        abstract* = {chapters/abstract-en.tex},
        keywords* = {voice cloning detection, speech anti-spoofing, AASIST, graph attention network, deepfake speech},
        acknowledgements = {chapters/acknowledgements.tex},
        bib-resource = {reference.bib}
    }
}

\begin{document}

\chapter{绪论}

\section{研究背景与意义}

近年来，随着深度学习技术的持续演进，文本转语音、语音转换以及个性化语音合成等技术取得了显著进展。以 WaveNet、Tacotron、FastSpeech、VITS 等为代表的语音生成模型在自然度、说话人相似性和语义可控性方面不断提升，使得语音克隆技术在数字人交互、智能客服、影视配音、无障碍服务等领域展现出较高的应用价值。然而，相关技术的快速发展也显著降低了语音伪造的门槛，攻击者可以借助公开语音片段生成高仿真的目标说话人语音，从而对远程身份认证、金融反欺诈、舆情传播和媒体内容可信性带来现实挑战\cite{wuzetal2015,kinnunen2017,todisco2019}。

传统自动说话人验证系统更关注说话人身份特征建模，对于语音内容是否由机器合成、转换或重放生成的判别能力相对有限。因此，构建能够有效识别伪造语音的反欺骗系统，已经成为语音安全研究中的重要课题。语音克隆检测的研究不仅具有突出的理论意义，也具有较强的工程应用价值：在理论层面，该任务涉及语音信号处理、深度学习、多域特征融合与图神经网络等交叉方向；在应用层面，该任务可服务于语音身份核验、安全审计、媒体审核与数字内容取证等实际场景。

\section{国内外研究现状}

语音反欺骗研究最初主要围绕重放攻击检测展开，随后逐步扩展至语音合成与语音转换攻击检测。ASVspoof 系列挑战赛的提出为该领域提供了标准化数据集与统一评测协议，有力推动了语音伪造检测方法的发展\cite{wuzetal2015,kinnunen2017,todisco2019,wang2020asvspoof,yamagishi2021}。早期研究通常依赖 MFCC、CQCC、LFCC 等人工设计特征，并结合高斯混合模型、支持向量机等传统分类器实现判别。这类方法在特定数据集上具有一定效果，但对未知攻击类型和跨域场景的泛化能力相对有限。

随着深度学习的发展，卷积神经网络、残差网络、时序卷积网络以及 Transformer 等方法逐渐成为主流。研究者开始尝试从梅尔频谱图、对数功率谱乃至原始波形中自动学习判别性表示，以提高模型对复杂伪造痕迹的建模能力\cite{jung2022aasist,tak2021rawnet2,velickovic2018}。其中，AASIST 模型通过联合建模谱域与时域信息，并引入异构图注意力机制，有效提升了单模型在语音反欺骗任务中的检测性能，已成为近年来具有代表性的研究成果之一\cite{jung2022aasist,tak2021stgat}。

\section{研究内容}

结合本毕业设计项目，本文围绕中文语音克隆检测系统的设计与实现开展研究，主要内容包括：

\begin{enumerate}
    \item 梳理语音克隆检测的研究背景、任务特点、常见攻击类型以及国内外研究进展；
    \item 分析语音预处理、重采样、定长裁剪及梅尔频谱特征提取等关键技术；
    \item 结合项目代码实现，对 AASIST 模型的整体架构与核心模块进行系统分析；
    \item 基于 FastAPI、PyTorch 与前端页面技术实现语音克隆检测软件系统，完成音频上传、格式转换、批量推理和结果展示等功能；
    \item 对系统的应用价值、局限性及后续改进方向进行总结与展望。
\end{enumerate}

\section{论文组织结构}

全文共分为五章，各章安排如下：第一章为绪论，介绍研究背景、研究现状、研究内容及论文结构；第二章阐述语音克隆检测相关技术，包括任务概述、基本原理、语音预处理与梅尔频谱特征；第三章对 AASIST 模型架构及其工程实现进行详细分析；第四章介绍语音克隆检测系统的设计与实现过程；第五章对全文工作进行总结，并对未来研究方向进行展望。

\chapter{声音克隆检测相关技术}

\section{声音克隆检测任务概述}

声音克隆检测又称语音伪造检测或语音反欺骗，其目标是在输入语音中识别是否存在由文本转语音、语音转换、重放或编辑拼接等人工过程所引入的异常痕迹。与自动说话人识别任务强调“是否为某位目标说话人”不同，语音克隆检测更关注“该语音是否由自然发声机制产生”。因此，检测模型需要对频带分布、相位连续性、韵律结构以及局部失真伪迹具备较强的辨识能力。

\section{克隆语音识别基本原理}

克隆语音识别的核心在于对真实语音与伪造语音之间的统计差异进行建模。真实语音由人类发声器官产生，通常具有连续且稳定的声学规律；而伪造语音虽然在主观听感上可能已较接近真实语音，但在局部时间结构、频谱细节以及跨频带耦合关系等方面仍可能存在可被模型学习的偏差。典型的检测流程包括输入音频预处理、特征表示构建、分类模型判别以及阈值决策输出等步骤。

\section{语音信号预处理方法}

在语音克隆检测任务中，预处理是保障模型输入一致性与稳定性的基础环节。常见处理步骤包括重采样、单声道化、幅值归一化、定长裁剪与填充等。本项目后端通过 `librosa` 将输入统一转换为 16 kHz 单声道语音，并使用定长化策略保证模型输入长度一致；对于多种上传音频格式，系统借助 `pydub` 与 `ffmpeg` 完成统一转换，以提升工程可用性。

\section{梅尔频谱特征提取流程}

梅尔频谱是语音分析中常用的时频特征，其提取过程通常包括分帧、加窗、短时傅里叶变换、梅尔滤波器组映射以及对数能量计算等步骤。该特征能够较好地保留语音的共振峰结构、谐波分布和能量变化规律，因此在传统语音识别和伪造语音检测任务中得到了广泛应用。尽管本文部署的核心模型采用原始波形输入，但对梅尔频谱提取机制的理解仍有助于把握语音反欺骗领域的基础理论。

\section{AASIST 模型简介}

AASIST 是一种面向语音反欺骗任务的端到端检测模型，其核心思想是联合建模语音信号中的谱域信息和时域信息，并通过图注意力机制挖掘节点之间的非局部依赖关系。该模型首先从原始波形中提取高层表示，然后分别构建谱域图和时域图，进一步通过异构图注意力层实现跨域信息融合，从而增强对伪造语音细粒度异常特征的刻画能力\cite{jung2022aasist,wang2019han}。

\chapter{AASIST 模型架构与实现}

\section{模型总体架构}

结合项目代码实现，AASIST 模型整体由时间卷积前端、残差编码器、谱域图注意力分支、时域图注意力分支、异构堆叠图注意力层、图池化模块以及读出分类层构成。模型输入为定长原始语音序列，输出为针对真实语音与伪造语音的二分类结果。项目中的关键参数配置存储于 `config_standalone_eval.json`，推理实现则封装于 `predictor.py`。

\section{输入层与时间卷积前端}

模型前向传播由原始波形输入开始。输入张量在增加通道维度后进入时间卷积前端，该模块以带通滤波器组思想为基础，通过可学习卷积核直接从波形中提取局部频带响应。相较于预先固定的特征变换方法，时间卷积前端能够在保留感知先验的同时，以数据驱动方式学习更适合反欺骗任务的特征表示。

\section{残差编码器与图注意力建模}

在前端卷积之后，模型利用多个残差块对高层时频表示进行深层抽象，以提升特征的判别能力与表达层次。随后，模型分别从编码器输出中构建谱域节点表示和时域节点表示，并通过图注意力机制学习节点之间的依赖关系。谱域分支主要关注不同频带之间的关系建模，时域分支则强调时间结构与韵律异常的判别。为进一步融合不同域的信息，AASIST 引入异构图注意力层与主节点更新机制，在统一图结构中实现时域与谱域信息的联合推理\cite{jung2022aasist,velickovic2018,wang2019han}。

\section{图池化与分类输出}

为降低图结构冗余并强化关键节点表达，模型在多个阶段使用图池化模块对节点进行筛选与加权。经过双分支异构图推理后，模型分别对时域表示和谱域表示执行最大池化与平均池化，并与主节点表示拼接形成最终隐藏向量，再通过线性分类层输出二分类 logits。项目推理流程以模型输出分数和预设阈值为依据，完成“真实”与“伪造”的最终判定。

\section{项目中的推理实现}

在工程实现中，`predictor.py` 对 AASIST 模型进行了完整封装。系统初始化时完成配置文件加载、设备选择、模型构建、权重加载与评估模式切换；推理时则依次执行音频读取、重采样、定长处理与批量张量构造，并调用模型完成统一推理。项目同时支持单文件预测和批量预测两类接口，从而较好地适配前端批量上传的业务需求。

\chapter{软件系统设计与实现}

\section{系统总体设计}

本项目实现了一个面向语音克隆检测场景的轻量级 Web 系统，整体由前端交互层、后端服务层和模型推理层组成。前端负责文件上传、试听与结果展示；后端负责接口接收、临时文件管理、音频格式转换与结果封装；模型推理层负责加载训练好的 AASIST 模型并输出检测结果。系统整体结构清晰，能够较好体现“算法分析—模型部署—界面交互”的完整设计思路。

\section{后端服务设计}

后端采用 FastAPI 框架实现，入口文件为 `app.py`。系统在启动阶段完成模型、配置文件及权重文件的检查，并将模型常驻内存以减少重复加载带来的时间开销。在请求处理过程中，后端通过 `POST /predict/` 接口接收多文件上传数据，对非 WAV 格式音频进行统一转换，随后调用批量推理接口完成检测，并将文件名、标签、分数及异常信息以结构化结果返回前端。

\section{前端界面设计}

前端页面由 `index.html` 实现，采用原生 HTML、CSS 与 JavaScript 完成界面布局与交互逻辑。页面支持点击上传和拖拽上传两种方式，并可对已选择音频进行本地试听。用户触发检测后，前端将所有文件封装为 `FormData` 并提交至后端接口，随后根据返回结果动态更新状态标签、分数信息和异常提示，从而形成完整的人机交互闭环。

\section{前后端联动过程}

系统运行流程可概括为：用户访问前端页面后上传音频文件，浏览器将文件信息与试听控件渲染至结果表格；用户点击检测按钮后，前端发起批量检测请求；后端接收文件并完成保存、格式转换、预处理与批量推理；模型输出结果后，后端将判别标签与分数返回前端，最终由前端完成结果可视化展示。该流程表明系统已具备较为完整的业务链路与基本工程落地能力。

\section{系统可扩展性分析}

尽管当前系统已实现基本的语音克隆检测功能，但仍具有进一步扩展空间。例如，可引入数据库对历史检测记录进行管理，增加健康检查接口以支持服务状态监控，并通过模型压缩、推理加速或前后端分离部署进一步提升系统的可用性与部署灵活性。这些方向均为后续工程优化提供了可行基础。

\chapter{总结与展望}

\section{全文总结}

本文围绕中文语音克隆检测任务，完成了从相关技术梳理、核心模型分析到软件系统实现的整体研究工作。论文系统分析了语音反欺骗的研究背景、任务特点与关键技术，在此基础上结合项目代码对 AASIST 模型的结构与实现机制进行了较为深入的讨论；同时，依托 FastAPI、PyTorch 与前端页面技术，设计并实现了一个可运行的语音克隆检测系统，完成了音频上传、格式转换、批量推理与结果展示等核心功能。整体来看，本文工作较好体现了毕业设计在理论研究、工程实现和应用场景结合方面的综合性。

\section{未来展望}

未来研究仍可从多个方面继续深化。首先，可围绕未知攻击类型检测能力开展研究，尝试引入多模型融合、特征蒸馏与阈值自适应策略，以提升模型泛化性能。其次，可结合更贴近真实业务场景的中文数据集，对模型在跨说话人、跨设备和跨平台环境下的稳定性进行系统评估。再次，随着语音生成大模型的发展，后续工作还可进一步探索开放集识别、攻击溯源与可解释分析等问题。最后，在系统部署层面，可通过 ONNX、TensorRT 或容器化方案提升推理效率与工程可维护性，以适应更加复杂的应用需求。

\nocite{*}

\end{document}
